
\section{Wissenschaftliche Herangehensweise}
\label{sec:method}

Zur Beantwortung der Forschungsfrage wird ein \textbf{Mixed-Methods-Ansatz} gewählt, der eine quantitative technische Analyse mit einer qualitativen Evaluation verbindet.

\subsection{Untersuchungsdesign \& Prototyp}
Es wird ein prototypischer Workflow bzw. eine Benutzeroberfläche entwickelt, die IDS-Informationen nutzt, um Anwender:innen beim Ergänzen von Daten zu leiten. Die unabhängige Variable ist dabei der Einsatz dieser IDS-basierten Benutzeroberfläche.

\subsection{Datenerhebung \& Stichprobe}
Zwei Datentypen werden erhoben:

\begin{enumerate}
    \item \textbf{Quantitative Analyse (Technische Validierung):}
    \begin{itemize}
        \item \textit{Material:} 5--10 repräsentative IFC-Modelle aus realen oder anonymisierten Projekten der Angebotskalkulation, die typische Defizite aufweisen.
        \item \textit{Methode:} Durchführung von IDS-Validierungen vor und nach der Bearbeitung mit dem Prototyp. Die Ergebnisse werden mithilfe des Open-Source-Tools \textit{IfcTester} dokumentiert und gegenübergestellt \cite{ifctester}.
    \end{itemize}
    \item \textbf{Qualitative Analyse (Expertenfeedback):}
    \begin{itemize}
        \item \textit{Stichprobe:} 2--3 Expert:inneninterviews mit Fachanwender:innen.
        \item \textit{Methode:} Nach der Testnutzung des Prototyps wird Feedback zur Praxistauglichkeit, Verständlichkeit und Zufriedenheit erhoben.
    \end{itemize}
\end{enumerate}

\subsection{Datenauswertung}
Die technischen Daten (Validierungsreports) werden vergleichend analysiert, um die Qualitätssteigerung der Modelle zu messen. Die qualitativen Interviews werden inhaltsanalytisch (in Anlehnung an Mayring) ausgewertet.
